\section{The angels in liturgy and devotion}

\subsection{The Roman celebrations}

The Roman Missal of 1962 celebrates the dedication of Saint Michael the
Archangel on 29 September as a first-class feast. Saint Gabriel has his own
third-class feast on 24 March, and Saint Raphael on 24 October. The Holy
Guardian Angels are celebrated on 2 October, also third class. These are the
titles and ranks of this edition; the dedication title on 29 September
belongs to the feast's liturgical identity.\footnote{\emph{Missale Romanum},
editio typica (Vatican, 1962), pp.~493, 670, 672, 695; CMAA facsimile,
PDF pp.~574, 751, 753, 776,
\url{https://media.churchmusicassociation.org/pdf/missale62.pdf}.}

The postconciliar General Roman Calendar joins Michael, Gabriel, and Raphael
in the feast of 29 September and keeps the Guardian Angels on 2 October as
an obligatory memorial. The \emph{Ordo lectionum Missae}, second typical
edition of 1981, appoints John 1:47--51 for the archangels: the angels ascend
and descend upon the Son of Man. Its first reading offers Daniel 7:9--10,
13--14 or Apocalypse 12:7--12a. For the Guardian Angels, Matthew 18:1--5, 10
is expressly the memorial's proper Gospel.\footnote{\emph{Ordo lectionum
Missae}, editio typica altera (Vatican, 1981), nn.~647 and 650, pp.~308--309,
PDF pp.~362--363; Congregation for Divine Worship and the Discipline of the
Sacraments, \emph{Directory on Popular Piety and the Liturgy}, n.~215,
for the postconciliar Roman celebrations.}

These selections place angelology within the history of salvation. The
Gospel of the archangels directs attention toward Christ, the meeting of
heaven and earth. The Guardian Angels' Gospel joins heavenly dignity to
care for the little ones. The person who can be overlooked on earth is
not negligible before God. Angelic ministry therefore deepens the Church's
attention to human persons rather than competing with it.

The 1962 formularies give the three named archangels distinct scriptural
settings. Gabriel's lesson is Daniel 9:21--26; Raphael's is Tobias
12:7--15. Michael's is Apocalypse 1:1--5. The common introit drawn from
Psalm 102:20 [Hebrew 103:20] places their particular missions within the
obedience and praise of all the heavenly ministers. A feast's focus on one
angel does not detach him from the common worship of God.\footnote{\emph{Missale
Romanum} (1962), pp.~493, 670, 695. Biblical assignments and themes are
reported here; the missal's liturgical texts are not reproduced.}

\subsection{The Byzantine synaxis}

The Byzantine Catholic Metropolitan Cantor Institute of the Archeparchy of
Pittsburgh publishes a Menaion office for 8 November under the title of the
Synaxis of the Holy Archangel Michael and all the Angelic Powers. Its
calendar marks a feast with an All-night Vigil. The celebration assembles
the heavenly hosts around the same ministry of praise and protection that
the name \emph{bodiless powers} expresses. This is the calendar and usage
witnessed by that particular Byzantine Catholic source.\footnote{Metropolitan
Cantor Institute, Archeparchy of Pittsburgh, online \emph{Menaion},
8 November, heading and feast symbol, accessed 29 September 2026,
\url{https://mci.archpitt.org/menaion/11-08.html}. No service text is reproduced.}

The Latin ranks first class, third class, feast, and memorial do not supply
a scale for this Byzantine celebration. Each belongs to its own liturgical
ordering. The common theological subject is nevertheless clear: the
heavenly ministers stand before God and serve his people. The celebration
is corporate as well as personal. Michael's prominence gathers attention
to the angelic host rather than making that host disappear behind a single
protecting figure.

\subsection{Prayer within Christ's worship}

The Church's worship joins earthly voices to heavenly praise. The
\emph{Directory on Popular Piety and the Liturgy} recalls the Sanctus,
angelic intercession, protection through the night, and the commendation
of the dying. It describes angelic devotion as gratitude to God, reverence,
and confidence in providence.\footnote{Congregation for Divine Worship and
the Discipline of the Sacraments, \emph{Directory on Popular Piety and the
Liturgy}, nn.~213--216; decree dated 17 December 2001, official English
web witness,
\url{https://www.vatican.va/roman_curia/congregations/ccdds/documents/rc_con_ccdds_doc_20020513_vers-direttorio_en.html}.}

The connection between worship and ministry is especially strong in Thomas.
Angels do not have to relinquish their possession of God in order to help
human beings. Their service proceeds from the order of divine providence.
For the worshipper, this makes gratitude more exact: God is thanked for a
created friend and minister whose excellence belongs to God's gift. The
minister may be honored and his intercession sought without being treated
as the source of the grace requested.\footnote{\ST{I, q.~112, a.~1,
corpus and ad~3--4}; \url{https://www.newadvent.org/summa/1112.htm}.}

\subsection{Names and the discipline of devotion}

The Directory's n.~217 directs that assigning names to angels be discouraged,
except for Michael, Gabriel, and Raphael, whose names occur in Scripture.
It also rejects an account of the world that leaves people helpless between
rival spiritual powers, and the habit of assigning every setback to a devil
and every success to a guardian angel. Its positive alternative is moral
commitment to the Gospel, humility, and prayer.\footnote{\emph{Directory},
n.~217. The wording is a curial pastoral directive about devotion, not a
definition of the number of angels who possess personal identity.}

The restriction gives devotion a determinate scriptural vocabulary. A
Christian need not discover a guardian angel's private name to acknowledge
his care. Nor does the absence of a supplied name make that ministry
impersonal. In the theological account, an angel knows and wills; in the
devotional act, the faithful person entrusts himself to the providence that
sends this minister. Knowledge of an undisclosed name contributes nothing
necessary to that relation.

Gratitude for protection also leaves room for suffering. The good that
providence seeks is the person's journey to God; temporal ease is not its
complete measure. A feast of the angels consequently teaches confidence
without supplying a method for interpreting every accident. The angels'
obedience directs the Christian toward his own obedience, and their praise
invites him into the common worship whose center is Christ.
